\section{An order-nine determinant correction at the same conductor}
\label{sec:cubic-determinant-ninth-root}

\begin{proposition}[AC-G46: an anticyclotomic cube root and determinant repair]
\label{AC-G46}
Let $K=\mathbb Q(\omega)$, let $\mathfrak q=(269)$ in
$\mathcal O_K=\mathbb Z[\omega]$, and let
$\chi=\chi_{269}$ be either choice of the nontrivial cubic
character for $K(\sqrt[3]{269})/K$, with the reciprocity
convention fixed as in AC-K04.  There is a finite character
\[
 \nu:G_K\longrightarrow\mu_9,
 \qquad \operatorname{ord}(\nu)=9,
 \qquad \nu^3=\chi,\qquad \nu^c=\nu^{-1},             \tag{G46.1}
\]
whose exact conductor is $\mathfrak q$.  It is trivial on the
entire local Galois group at the prime above $3$.  At every
rational prime $p\ne269$ inert in $K$, it is unramified and
\[
                     \nu(\operatorname{Frob}_{(p)})=1. \tag{G46.2}
\]

Consequently, for any three-dimensional representation
$\rho:G_K\to\operatorname{GL}_3(L)$ with
$\det\rho=\chi^2$, where the coefficient field contains
$\mu_9$, the twist $\rho^{\mathrm{det}}=\nu\rho$ satisfies
\[
       \det\rho^{\mathrm{det}}=1,
       \operatorname{End}(\rho^{\mathrm{det}})
                       =\operatorname{End}(\rho)        \tag{G46.3}
\]
as representations in the chosen matrix realization.
The same equality holds for the source semidirect adjoint
representation when a normalized lift to $\mathcal G_3$ is
present.  The twist preserves irreducibility and Hodge--Tate
weights, is unchanged locally above $3$, and leaves the full
Frobenius matrix unchanged at the primes in (G46.2).

For the proposed cubic diagonal seed it gives
\[
 \nu\operatorname{diag}(\chi\gamma,1,\chi\gamma^{-1})
     =\operatorname{diag}(\nu^4\gamma,\nu,
                          \nu^4\gamma^{-1}),            \tag{G46.4}
\]
whose determinant is $\nu^9=1$.
If its prescribed inertia type at $\mathfrak q$ is the tame
type $\chi\oplus1\oplus\chi$, the original conductor
exponent there is $2$, and the corrected exponent is $3$.
At other finite primes the conductor exponents do not change.
This conductor assertion uses the usual Weil--Deligne conductor
at a coefficient prime distinct from $269$.

This removes determinant one as an algebraic obstruction to
this scalar correction.  It does not turn (G46.4) into the
source seed $\operatorname{diag}(\Psi,1,\Psi^{-1})$ or
transport the source \texttt{Dat}, automorphy, rigidity, or
cohomology extraction theorem.
\end{proposition}
\begin{proof}
\proofstep{1. The exact ray group.}
The prime $269$ is $2$ modulo $3$, hence is inert in $K$ and
$\mathcal O_K/\mathfrak q\simeq\mathbb F_{269^2}$.
The class number is one and the unit group is $\mu_6$.
The finite ray class exact sequence therefore gives
\[
 \operatorname{Cl}_{\mathfrak q}(K)
       \simeq\mathbb F_{269^2}^{\times}/\mu_6,
 \qquad
 \#\operatorname{Cl}_{\mathfrak q}(K)
       =\frac{269^2-1}{6}
       =12060=2^2\cdot3^2\cdot5\cdot67.                \tag{G46.5}
\]
The reduction of $\mu_6$ is injective: a nontrivial sixth root
can become $1$ only in residue characteristic $2$ or $3$.
The multiplicative group of a finite field is cyclic, so the
quotient in (G46.5) is cyclic.  These are the same ray modulus,
unit quotient, and fixed cubic character used in AC-K04 and
AC-G40; no larger modulus is introduced here.

The subgroup of characters with values in $\mu_9$ is cyclic
of order $9$.  Its cube map has as image the whole group of
cubic characters.  Choose a lift $\nu$ of the fixed $\chi$.
Since $\nu^3$ has order $3$, $\nu$ has order exactly $9$.
The character has conductor dividing $\mathfrak q$; the only
proper divisor is the unit modulus.  An everywhere unramified
abelian character of $K$ factors through its trivial ideal class
group, whereas $\nu$ is nontrivial.  Its conductor is therefore
exactly $\mathfrak q$.

\proofstep{2. Conjugation and the local trivialities.}
On $\mathbb F_{269^2}$ complex conjugation induces the
nontrivial residue-field automorphism, $x\mapsto x^{269}$.
Hence for every ray class $a$,
\[
          \nu(ca)=\nu(a)^{269}=\nu(a)^{-1},            \tag{G46.6}
\]
because $269\equiv-1\pmod9$.  Compatibility of reciprocity
with conjugation gives $\nu^c=\nu^{-1}$, with the same
conclusion under either consistent Artin convention.

Let $\lambda=1-\omega$ and $\mathfrak l=(\lambda)$.
The character is unramified at $\mathfrak l$ because its
conductor is $\mathfrak q$.  The equality
$\overline\lambda=-\omega^2\lambda$ shows that conjugation
fixes the ray class of $\mathfrak l$ (indeed it fixes this
prime ideal itself).  Equation (G46.6) gives
$\nu(\operatorname{Frob}_{\mathfrak l})^2=1$.
Its value has odd order dividing $9$, so it is $1$.
An unramified local character is determined by its Frobenius
value; thus the whole local character at $3$ is trivial.

For an inert rational prime $p\ne269$, the unique prime
of $K$ above it is the principal ideal $(p)$.
Its ray class is represented, up to the fixed reciprocity
inverse, by the reduction of the rational integer $p$ in
$\mathbb F_{269}^{\times}\subset\mathbb F_{269^2}^{\times}$.
This subgroup has order $268$, which is coprime to $9$.
A character of order dividing $9$ is therefore trivial on
it, proving (G46.2).

\proofstep{3. The representation and the exact scope.}
Scalar multiplication by $\nu$ changes the determinant by
$\nu^3=\chi$.  Therefore
$\det(\nu\rho)=\chi\chi^2=1$.
For every endomorphism $X$,
\[
       (\nu(g)\rho(g))X(\nu(g)\rho(g))^{-1}
                    =\rho(g)X\rho(g)^{-1}.             \tag{G46.7}
\]
For a normalized semidirect lift, extend $\nu$ to the
crossed character $a(hc^u)=\nu(h)$ as in (G45.4), and
multiply its matrix component by $a(g)$.
The scalar is inverted by $\vartheta$, so this again defines
a homomorphism and leaves the matrix at $c$ equal to $I$.
The source formula for \texttt{GG.adFun} places that scalar
and its inverse on the two sides of either $X$ or
$\texttt{thetaLie}(X)$; they cancel exactly.  Thus its
entire adjoint representation is unchanged, including the
mixed coefficient (G45.6).

Finite characters have Hodge--Tate weight zero.  Tensoring by
a nonzero scalar character preserves invariant subspaces,
and (G46.2) makes the Frobenius scalar equal to $1$ at the
specified primes.  Local triviality at $3$ was proved above.
At $\mathfrak q$ the two old nontrivial inertia characters
are $\chi=\nu^3$, and after correction the three characters
are $\nu^4,\nu,\nu^4$.  Each has order $9$, so all are
nontrivial; all are tame since the residue characteristic is
$269$.  The Artin conductor of a tame representation is its
dimension minus the dimension of its inertia invariants.
This gives respectively $3-1=2$ and $3-0=3$.
At every other prime $\nu$ is unramified, so the inertia
representation and Weil--Deligne monodromy, hence its conductor,
are unchanged.

The middle character in (G46.4) is $\nu$, not $1$.
In particular this corrected integral datum still is not the
literal source \texttt{r0} or \texttt{Dat}; at a coefficient
prime different from $3$, its nontrivial middle character is
also visible after reduction.  At the prime $3$ its reduction
may be trivial, but its characteristic-zero middle character
and the defining datum remain different.  The old universal
quantifiers are therefore not supplied by the determinant
calculation.  Finally (G46.7) shows directly why this repair
does not prove the new Type II vanishing (G45.7).

This distinction is also literal in the counting interface:
\path{Development/AutomorphicArithmetic/Foundations/GaloisRepresentations/AdmissibleResidues/IsAutomorphicOverImagQuad.lean},
\texttt{AutCsd0}, lines 35--36, requires a regular algebraic
cuspidal representation, conjugate duality, and attachment.
It does not require trivial central character.  Likewise
\texttt{Adm3Alt}, lines 29--33 of
\path{Development/AutomorphicArithmetic/AutomorphicVarieties/ModularityArgument/ThresholdArguments/DenseThresholdCases/AfStmt.lean},
and \texttt{Adm3U\_s} in that directory's
\path{NumberU3US.lean}, require irreducibility, the stated
automorphic witness, level, and unlinking data, without an
explicit determinant-one clause.  Thus (G46.3) is an available
normalization; it is not being credited with proving an extra
determinant hypothesis of that generic counting theorem.
\end{proof}

\paragraph{Geometric dictionary.}
The kernel of $\nu$ defines a cyclic degree-nine extension of
$K$ containing the cubic extension cut out by $\chi=\nu^3$.
Its integral normalization is a finite \'etale cover of
$\operatorname{Spec}\mathcal O_K[1/(3\cdot269)]$.
In the already audited SGA~1, Expos\'e~V, \S\S4 and~7
dictionary, the map from the degree-nine cover to its degree-three
quotient is precisely restriction along the quotient of the
cyclic Galois groups.  Its cubic quotient has the Kummer
interpretation of SGA~1, Expos\'e~XI, \S6,
Propositions~1 and~4 and Corollary~5.
The degree-nine cover is asserted here as a cyclic ray cover;
no choice of a ninth-root Kummer equation over $K$ is assumed.
Only the previously audited framework locators are reused.
